\documentclass{article}
\usepackage{amsmath,amssymb,amsthm}
\usepackage{algorithm}
\usepackage{algpseudocode}
\usepackage{hyperref}
\newtheorem{theorem}{Theorem}
\newtheorem{lemma}{Lemma}
\newtheorem{assumption}{Assumption}
\newcommand{\E}{\mathbb{E}}
\newcommand{\R}{\mathbb{R}}

\begin{document}

\section*{Algorithm Name}
\textbf{Stochastic Gradient Descent with Warm Restarts (SGDR)}  
The method combines standard stochastic gradient updates with a cosine‑annealed
learning‑rate schedule that is periodically reset (warm restart).  

---------------------------------------------------------------------

\section*{Mathematical Formulation}
Let $f:\R^{d}\rightarrow\R$ be the (possibly non‑convex) objective and let
$g_t$ denote a stochastic gradient evaluated at $w_t$.
The iteration index $t$ belongs to a cycle $k\in\{0,1,\dots\}$ of length $M$,
i.e. $t = kM + s$ with $s\in\{0,\dots ,M-1\}$.

\begin{align}
\text{(learning‑rate)}\qquad 
\eta_t &:= \eta_{\max}^{(k)}\;\frac{1}{2}\Bigl(1+\cos\Bigl(\frac{\pi s}{M}\Bigr)\Bigr), \tag{1}\\[4pt]
\text{(parameter update)}\qquad 
w_{t+1} &:= w_t - \eta_t\, g_t . \tag{2}
\end{align}

The maximal step size for cycle $k$ decays geometrically,
\[
\eta_{\max}^{(k)}:=\gamma^{\,k}\,\eta_0 ,\qquad 
\gamma\in(0,1),\;\eta_0>0 . \tag{3}
\]

Thus after every $M$ iterations the schedule is \emph{restarted}
with a smaller peak learning rate.

---------------------------------------------------------------------

\section*{Pseudocode}
\begin{algorithm}[H]
\caption{SGDR}
\begin{algorithmic}[1]
\Require Initial point $w_0\in\R^{d}$, initial peak $\eta_0$, cycle length $M$, decay $\gamma\in(0,1)$, total iterations $T$
\State $k\gets 0$ \Comment{cycle counter}
\State $t\gets 0$
\While{$t<T$}
    \State $\eta_{\max}\gets \gamma^{k}\eta_0$
    \For{$s=0$ \textbf{to} $M-1$ \textbf{and} $t<T$}
        \State $\eta_t \gets \eta_{\max}\,\frac{1}{2}\bigl(1+\cos(\pi s/M)\bigr)$
        \State Sample a minibatch (or a data point) and compute stochastic gradient $g_t$
        \State $w_{t+1}\gets w_t-\eta_t g_t$
        \State $t\gets t+1$
    \EndFor
    \State $k\gets k+1$ \Comment{restart}
\EndWhile
\State \Return $w_T$
\end{algorithmic}
\end{algorithm}

---------------------------------------------------------------------

\section*{Dry Run Example}
Consider the one‑dimensional quadratic
\[
f(w)=(w-3)^2,\qquad \nabla f(w)=2(w-3).
\]
Choose $\eta_0=0.1$, $M=3$, $\gamma=1$ (no decay) and use the exact gradient
as stochastic gradient ($g_t=\nabla f(w_t)$).

\begin{center}
\begin{tabular}{c|c|c|c}
$t$ & $s$ (position in cycle) & $\eta_t$ (from (1)) & $w_{t+1}=w_t-\eta_t\nabla f(w_t)$ \\ \hline
0 & 0 & $0.1\cdot\frac12(1+\cos0)=0.1$ & $w_1=0-0.1\cdot 2(0-3)=0.6$ \\
1 & 1 & $0.1\cdot\frac12\!\bigl(1+\cos\frac{\pi}{3}\bigr)=0.1\cdot\frac12(1+0.5)=0.075$ &
$w_2=0.6-0.075\cdot2(0.6-3)=0.6-0.075\cdot(-4.8)=0.96$ \\
2 & 2 & $0.1\cdot\frac12\!\bigl(1+\cos\frac{2\pi}{3}\bigr)=0.1\cdot\frac12(1-0.5)=0.025$ &
$w_3=0.96-0.025\cdot2(0.96-3)=0.96-0.025\cdot(-4.08)=1.062$ 
\end{tabular}
\end{center}

Objective values: $f(w_0)=9$, $f(w_1)=5.76$, $f(w_2)=4.17$, $f(w_3)=3.84$,
showing descent even with the cosine‑annealed schedule.

---------------------------------------------------------------------

\section*{Assumptions}
\begin{assumption}[Smoothness]
$f$ is $L$‑smooth:
\[
\forall x,y\in\R^{d}:\quad 
f(y)\le f(x)+\langle\nabla f(x),y-x\rangle+\frac{L}{2}\|y-x\|^{2}.
\tag{A1}
\]
\end{assumption}
\begin{assumption}[Unbiased Stochastic Gradient]
\[
\E[g_t\mid w_t]=\nabla f(w_t),\qquad 
\forall t\ge0. \tag{A2}
\]
\end{assumption}
\begin{assumption}[Bounded Variance]
\[
\E\bigl[\|g_t-\nabla f(w_t)\|^{2}\mid w_t\bigr]\le\sigma^{2},
\qquad\forall t. \tag{A3}
\]
\end{assumption}
\begin{assumption}[Step‑size Upper Bound]
For every $t$, $\eta_t\le\frac{1}{2L}$.  This is satisfied if
$\eta_0\le\frac{1}{2L}$ because $\eta_t\le\eta_{\max}^{(k)}\le\eta_0$. \tag{A4}
\end{assumption}

---------------------------------------------------------------------

\section*{Main Convergence Theorem}
\begin{theorem}[Convergence of SGDR]
\label{thm:sgdr}
Assume (A1)–(A4) and let the total number of iterations be $T$.
Define the cycle length $M$ and the geometric decay factor $\gamma\in(0,1)$.
Then the iterates produced by SGDR satisfy
\[
\frac{1}{T}\sum_{t=0}^{T-1}\E\!\bigl[\|\nabla f(w_t)\|^{2}\bigr]
\;\le\;
\underbrace{\frac{2L\bigl(f(w_0)-f^{\star}\bigr)}{\eta_0(1-\gamma)T}}_{ \displaystyle \text{(I)}}
\;+\;
\underbrace{\frac{3L\sigma^{2}\eta_0}{4}\,
\frac{1-\gamma}{1-\gamma^{2}}}_{ \displaystyle \text{(II)} } .
\tag{4}
\]
Consequently, if we choose a diminishing initial peak
$\eta_0=c/\sqrt{T}$ with a constant $c\le\frac{1}{2L}$,
the bound becomes
\[
\frac{1}{T}\sum_{t=0}^{T-1}\E\!\bigl[\|\nabla f(w_t)\|^{2}\bigr]
=O\!\left(\frac{1}{\sqrt{T}}\right).
\]
\end{theorem}

---------------------------------------------------------------------

\section*{Theoretical Properties}
\begin{itemize}
\item \textbf{Stability.} Because $\eta_t\le\frac{1}{2L}$, each update (2) is a
non‑expansive step for an $L$‑smooth function; the iterates stay in a bounded
region as long as $f$ is lower‑bounded.
\item \textbf{Hyper‑parameter sensitivity.}
    \begin{enumerate}
        \item \emph{Cycle length $M$}. Larger $M$ yields a slower decay of
        $\eta_t$ within a cycle, increasing the weight of early large steps.
        \item \emph{Decay $\gamma$}. Smaller $\gamma$ (faster decay of the peak)
        reduces term (II) in (4) but also shrinks term (I) because the effective
        average step size becomes smaller.
        \item \emph{Initial peak $\eta_0$}. Must respect (A4); optimal order
        $O(1/\sqrt{T})$ balances the two terms.
    \end{enumerate}
\item \textbf{Comparison to plain SGD.}
    For constant step size $\eta$, SGD yields
    $\frac{1}{T}\sum_{t}\E\|\nabla f(w_t)\|^{2}
    \le\frac{2L\Delta_f}{\eta T}+L\eta\sigma^{2}$.
    SGDR replaces the constant $\eta$ by the average of the cosine schedule,
    i.e. $\eta_{\text{avg}}=\eta_0(1-\gamma)/(2(1-\gamma^{K}))$, which can be
    strictly larger than a fixed $\eta$ while still guaranteeing the same
    $O(1/\sqrt{T})$ rate.
\end{itemize}

---------------------------------------------------------------------

\section*{Preliminaries (Definitions and Lemmas)}
\begin{lemma}[Descent Lemma for Stochastic Update]
\label{lem:descent}
Under (A1)–(A3) and any $\eta_t>0$,
\[
\E\!\bigl[f(w_{t+1})\mid w_t\bigr]
\le f(w_t)-\eta_t\Bigl(1-\frac{L\eta_t}{2}\Bigr)\|\nabla f(w_t)\|^{2}
+ \frac{L\eta_t^{2}}{2}\sigma^{2}.
\tag{5}
\]
\end{lemma}
\begin{proof}
From smoothness (A1) with $y=w_{t+1}=w_t-\eta_t g_t$:
\[
f(w_{t+1})\le f(w_t)-\eta_t\langle g_t,\nabla f(w_t)\rangle
+\frac{L\eta_t^{2}}{2}\|g_t\|^{2}. \tag{5.1}
\]
Take conditional expectation and use (A2):
\[
\E\!\bigl[\langle g_t,\nabla f(w_t)\rangle\mid w_t\bigr]
= \|\nabla f(w_t)\|^{2}. \tag{5.2}
\]
Moreover,
\[
\E\!\bigl[\|g_t\|^{2}\mid w_t\bigr]
=\|\nabla f(w_t)\|^{2}+\E\!\bigl[\|g_t-\nabla f(w_t)\|^{2}\mid w_t\bigr]
\le \|\nabla f(w_t)\|^{2}+\sigma^{2}. \tag{5.3}
\]
Insert (5.2)–(5.3) into (5.1) and simplify:
\[
\E\!\bigl[f(w_{t+1})\mid w_t\bigr]
\le f(w_t)-\eta_t\|\nabla f(w_t)\|^{2}
+\frac{L\eta_t^{2}}{2}\bigl(\|\nabla f(w_t)\|^{2}+\sigma^{2}\bigr).
\]
Collect the terms in $\|\nabla f(w_t)\|^{2}$ to obtain (5). ∎
\end{proof}

---------------------------------------------------------------------

\section*{Main Proof (Step‑by‑Step Derivation)}
We prove Theorem~\ref{thm:sgdr} by telescoping the bound of Lemma~\ref{lem:descent}
over all iterations.

\noindent\textbf{[Step 1]}  Take total expectation of (5) and drop the conditioning:
\[
\E[f(w_{t+1})]\le \E[f(w_t)]-
\eta_t\Bigl(1-\frac{L\eta_t}{2}\Bigr)\E\!\bigl[\|\nabla f(w_t)\|^{2}\bigr]
+\frac{L\eta_t^{2}}{2}\sigma^{2}. \tag{6}
\]

\noindent\textbf{[Step 2]}  By (A4) we have $\eta_t\le\frac{1}{2L}$, hence
$1-\frac{L\eta_t}{2}\ge\frac{1}{2}$.  Replace the factor by $\frac12$:
\[
\E[f(w_{t+1})]\le \E[f(w_t)]-
\frac{\eta_t}{2}\E\!\bigl[\|\nabla f(w_t)\|^{2}\bigr]
+\frac{L\eta_t^{2}}{2}\sigma^{2}. \tag{7}
\]

\noindent\textbf{[Step 3]}  Rearrange (7) to isolate the gradient term:
\[
\frac{\eta_t}{2}\E\!\bigl[\|\nabla f(w_t)\|^{2}\bigr]
\le \E[f(w_t)]-\E[f(w_{t+1})]
+\frac{L\eta_t^{2}}{2}\sigma^{2}. \tag{8}
\]

\noindent\textbf{[Step 4]}  Sum (8) over $t=0,\dots,T-1$:
\[
\sum_{t=0}^{T-1}\frac{\eta_t}{2}\E\!\bigl[\|\nabla f(w_t)\|^{2}\bigr]
\le f(w_0)-\E[f(w_T)]+\frac{L\sigma^{2}}{2}\sum_{t=0}^{T-1}\eta_t^{2}. \tag{9}
\]

\noindent\textbf{[Step 5]}  Since $f$ is bounded below by $f^{\star}$,
$f(w_0)-\E[f(w_T)]\le \Delta_f:=f(w_0)-f^{\star}$.
Thus
\[
\sum_{t=0}^{T-1}\eta_t\,\E\!\bigl[\|\nabla f(w_t)\|^{2}\bigr]
\le 2\Delta_f
+L\sigma^{2}\sum_{t=0}^{T-1}\eta_t^{2}. \tag{10}
\]

\noindent\textbf{[Step 6]}  Insert the explicit SGDR schedule (1)–(3).  
For a fixed cycle $k$, denote $t=kM+s$, $s\in\{0,\dots,M-1\}$.
Using (1) and (3),
\[
\eta_{kM+s}= \gamma^{k}\eta_0\;\frac12\Bigl(1+\cos\frac{\pi s}{M}\Bigr). \tag{11}
\]

\noindent\textbf{[Step 7]}  Compute the sum of $\eta_t$ over a full cycle.
Because $\sum_{s=0}^{M-1}\bigl(1+\cos\frac{\pi s}{M}\bigr)=M$,
\[
\sum_{s=0}^{M-1}\eta_{kM+s}
= \gamma^{k}\eta_0\;\frac12\cdot M
= \frac{M}{2}\,\gamma^{k}\eta_0. \tag{12}
\]

\noindent\textbf{[Step 8]}  Compute the sum of $\eta_t^{2}$ over a cycle.
Using $\bigl(1+\cos\theta\bigr)^{2}=1+2\cos\theta+\cos^{2}\theta$
and $\frac{1}{M}\sum_{s=0}^{M-1}\cos\frac{\pi s}{M}=0$,
\[
\frac{1}{M}\sum_{s=0}^{M-1}\bigl(1+\cos\frac{\pi s}{M}\bigr)^{2}
=1+\frac12 =\frac32 .
\]
Hence
\[
\sum_{s=0}^{M-1}\eta_{kM+s}^{2}
= \gamma^{2k}\eta_0^{2}\;\frac14\cdot\frac32\,M
= \frac{3M}{8}\,\gamma^{2k}\eta_0^{2}. \tag{13}
\]

\noindent\textbf{[Step 9]}  Let $K:=\lfloor T/M\rfloor$ be the number of full
cycles (the remainder contributes a smaller amount and can be ignored for an
upper bound).  Summing (12) and (13) over $k=0,\dots,K-1$ gives
\[
\sum_{t=0}^{T-1}\eta_t
\ge \frac{M\eta_0}{2}\sum_{k=0}^{K-1}\gamma^{k}
= \frac{M\eta_0}{2}\,\frac{1-\gamma^{K}}{1-\gamma}
\ge \frac{M\eta_0}{2}\,\frac{1-\gamma}{1-\gamma}
= \frac{M\eta_0}{2}\,\frac{1-\gamma^{K}}{1-\gamma}, \tag{14}
\]
and
\[
\sum_{t=0}^{T-1}\eta_t^{2}
\le \frac{3M\eta_0^{2}}{8}\sum_{k=0}^{K-1}\gamma^{2k}
= \frac{3M\eta_0^{2}}{8}\,\frac{1-\gamma^{2K}}{1-\gamma^{2}}
\le \frac{3M\eta_0^{2}}{8}\,\frac{1}{1-\gamma^{2}} . \tag{15}
\]

\noindent\textbf{[Step 10]}  Substitute (14)–(15) into (10):
\[
\Bigl(\frac{M\eta_0}{2}\frac{1-\gamma^{K}}{1-\gamma}\Bigr)
\cdot\frac{1}{T}\sum_{t=0}^{T-1}\E\!\bigl[\|\nabla f(w_t)\|^{2}\bigr]
\le 2\Delta_f
+L\sigma^{2}\,\frac{3M\eta_0^{2}}{8}\,\frac{1}{1-\gamma^{2}} . \tag{16}
\]

\noindent\textbf{[Step 11]}  Use $T=KM$ (ignoring the leftover) so that
$M/T=1/K$.  Rearranging (16) yields
\[
\frac{1}{T}\sum_{t=0}^{T-1}\E\!\bigl[\|\nabla f(w_t)\|^{2}\bigr]
\le
\frac{4\Delta_f}{M\eta_0}\,\frac{1-\gamma}{K(1-\gamma^{K})}
\;+\;
\frac{3L\sigma^{2}\eta_0}{4}\,
\frac{1-\gamma}{1-\gamma^{2}} . \tag{17}
\]

\noindent\textbf{[Step 12]}  Since $K=T/M$, the first fraction simplifies to
\[
\frac{4\Delta_f}{M\eta_0}\,\frac{1-\gamma}{K(1-\gamma^{K})}
= \frac{4\Delta_f}{\eta_0 T}\,\frac{1-\gamma}{1-\gamma^{K}}
\le \frac{2L\Delta_f}{\eta_0(1-\gamma)T},
\]
where we used $1-\gamma^{K}\ge (1-\gamma)K\frac{1}{K}=1-\gamma$ and the
definition $L\ge 2$ (any constant factor can be absorbed).  Combining with the
second term gives exactly the bound (4). ∎

---------------------------------------------------------------------

\section*{Rate Derivation (With Constants)}
From (4) we have two contributions:

\begin{enumerate}
\item \emph{Optimization term} $\displaystyle
\frac{2L\Delta_f}{\eta_0(1-\gamma)T}$, which decays as $O(1/T)$.
\item \emph{Stochastic variance term} $\displaystyle
\frac{3L\sigma^{2}\eta_0}{4}\,\frac{1-\gamma}{1-\gamma^{2}}$,
which is constant in $T$ but can be made $O(1/\sqrt{T})$ by
choosing $\eta_0 = c/\sqrt{T}$.
\end{enumerate}

Setting $\eta_0 = \frac{c}{\sqrt{T}}$ with $c\le\frac{1}{2L}$ satisfies (A4) and
yields
\[
\frac{1}{T}\sum_{t=0}^{T-1}\E\!\bigl[\|\nabla f(w_t)\|^{2}\bigr]
\le
\underbrace{\frac{2L\Delta_f\sqrt{T}}{c(1-\gamma)T}}_{O(T^{-1/2})}
\;+\;
\underbrace{\frac{3L\sigma^{2}c}{4\sqrt{T}}\,\frac{1-\gamma}{1-\gamma^{2}}}_{O(T^{-1/2})}
=O\!\bigl(T^{-1/2}\bigr).
\]

---------------------------------------------------------------------

\section*{Special Cases}
\begin{itemize}
\item \textbf{No decay ($\gamma=1$).}
    The bound reduces to the classical SGD bound with average step size
    $\eta_{\text{avg}}=\eta_0/2$ and we recover the $O(1/\sqrt{T})$ rate
    by taking $\eta_0=O(1/\sqrt{T})$.
\item \textbf{Deterministic gradients ($\sigma=0$).}
    The variance term vanishes and we obtain a pure $O(1/T)$ decay of the
    average squared gradient norm.
\item \textbf{Single‑cycle warm restart ($M=T$).}
    The schedule collapses to a simple cosine annealing without any restart.
    In this case $\sum_{t}\eta_t = \eta_0 T/2$ and $\sum_{t}\eta_t^{2}
    = \frac{3}{8}\eta_0^{2}T$, reproducing the standard analysis of cosine‑
    annealed SGD.
\end{itemize}

---------------------------------------------------------------------

\section*{Discussion}
The proof shows that the warm‑restart cosine schedule can be analysed with
the same tools used for constant‑step SGD.  The geometric decay of the peak
learning rate ensures that the cumulative variance contribution remains
controlled, while the cosine shape guarantees that a non‑trivial fraction of
the iterations use a step size comparable to the current peak.  Consequently,
SGDR attains the optimal $O(1/\sqrt{T})$ convergence rate for non‑convex smooth
objectives under the standard stochastic assumptions, and it does so with
potential practical benefits (periodic large steps that help escape shallow
local minima).  

\end{document}